\documentclass[12pt]{article}
\usepackage[utf8]{inputenc}
\usepackage[T1]{fontenc}
\usepackage[margin=1in]{geometry}
\usepackage{setspace}
\usepackage{xcolor}
\usepackage[colorlinks=true, urlcolor=blue!60!black, linkcolor=blue!60!black]{hyperref}
\setlength{\parindent}{2em}
\setlength{\parskip}{1.25em}
\renewcommand{\baselinestretch}{1.0}
\begin{document}
\begin{center}
{\LARGE \textit{Trump v. United States}, 603 U.S. 593 (2024)\par}
\vspace{0.8em}
{\large Case Brief\par}
\end{center}
\vspace{1em}

\section{Memory Jogger}

In \href{https://www.courtlistener.com/opinion/10600038/trump-v-united-states/}{\textit{Trump v. United States}}, the Supreme Court held that a former President enjoys absolute criminal immunity for exercises of core, exclusive constitutional powers, at least presumptive immunity for all other official acts within the outer perimeter of presidential responsibility, and no immunity for unofficial acts.

\section{Detailed Case Facts}

Donald J. Trump served as President of the United States from January 2017 until January 2021. Following his loss in the November 2020 presidential election to Joseph R. Biden, Jr., Special Counsel Jack Smith---appointed by Attorney General Merrick Garland in November 2022---led a federal grand jury investigation in the District of Columbia. On August 1, 2023, the grand jury returned a four-count indictment charging Trump with federal crimes arising out of conduct during his final months in office (\href{https://www.law.cornell.edu/supremecourt/text/23-939}{603 U.S. 593 (2024)}; \href{https://clearinghouse.net/case/46788/}{Civil Rights Litigation Clearinghouse}).

The indictment alleged that Trump conspired to overturn the legitimate results of the 2020 presidential election by spreading knowingly false claims of outcome-determinative election fraud to obstruct the collecting, counting, and certifying of the electoral votes. Specifically, the indictment alleged five primary means by which Trump and his co-conspirators pursued this objective (\href{https://www.law.cornell.edu/supremecourt/text/23-939}{Cornell LII}):
\begin{enumerate}
\item Using knowingly false claims of election fraud to pressure state legislators and election officials in Arizona, Georgia, Michigan, Pennsylvania, and Wisconsin to alter the popular vote count and replace legitimate electors with electors pledged to Trump;
\item Organizing fraudulent slates of electors in seven targeted states (Arizona, Georgia, Michigan, Nevada, New Mexico, Pennsylvania, and Wisconsin) and causing them to transmit false certificates to the Vice President, Congress, and the Archivist of the United States;
\item Attempting to leverage the Department of Justice to conduct sham election-fraud investigations and send an official letter to targeted states falsely asserting that the Department had identified significant concerns impacting the election outcome, while repeatedly threatening to replace the Acting Attorney General when he resisted;
\item Pressuring Vice President Mike Pence to use his ceremonial role presiding over the January 6, 2021, joint session of Congress to reject legitimate electoral votes or return them to state legislatures, and subsequently addressing gathered supporters on the morning of January 6 and directing them to the Capitol; and
\item Exploiting the violent attack on the Capitol by redoubling efforts to convince Members of Congress to delay certification of the electoral vote.
\end{enumerate}

Based on these allegations, the indictment charged Trump with four federal offenses: (1) conspiracy to defraud the United States under 18 U.S.C. \S 371; (2) conspiracy to obstruct an official proceeding under 18 U.S.C. \S 1512(k); (3) obstruction of and attempt to obstruct an official proceeding under 18 U.S.C. \S\S 1512(c)(2) and 2; and (4) conspiracy against rights under 18 U.S.C. \S 241 (\href{https://www.law.cornell.edu/supremecourt/text/23-939}{Cornell LII}).

\section{Procedural History}

Trump moved to dismiss the indictment in the United States District Court for the District of Columbia, asserting that a former President possesses absolute immunity from federal criminal prosecution for actions within the outer perimeter of his official responsibilities and that the Impeachment Judgment Clause (U.S. Const. art. I, \S 3, cl. 7) bars criminal prosecution unless the President is first impeached by the House and convicted by the Senate. On December 1, 2023, Judge Tanya S. Chutkan denied the motion to dismiss, holding that former Presidents enjoy no federal criminal immunity for acts committed while in office, without deciding whether the indicted conduct constituted official acts (\href{https://www.oyez.org/cases/2023/23-939}{Oyez}; \textit{United States v. Trump}, 704 F. Supp. 3d 196 (D.D.C. 2023)).

Special Counsel Smith petitioned the Supreme Court for certiorari before judgment, which the Court denied on December 22, 2023 (144 S. Ct. 539). On February 6, 2024, a unanimous three-judge panel of the U.S. Court of Appeals for the D.C. Circuit affirmed Judge Chutkan's denial of immunity in a per curiam opinion (\textit{United States v. Trump}, 91 F.4th 1173 (D.C. Cir. 2024)), reasoning that ``former President Trump has become citizen Trump'' and that separation-of-powers principles do not shield a former President from prosecution under generally applicable criminal laws (\href{https://www.law.cornell.edu/supremecourt/text/23-939}{Cornell LII}). Like the District Court, the D.C. Circuit declined to classify the indicted conduct as official or unofficial.

On February 28, 2024, the Supreme Court construed Trump's stay application as a petition for certiorari and granted review on the question whether, and to what extent, a former President enjoys presidential immunity from criminal prosecution for conduct alleged to involve official acts during his tenure in office (\href{https://www.supremecourt.gov/opinions/23pdf/23-939_e2pg.pdf}{Slip Op.}). The Court heard oral argument on April 25, 2024, and issued its decision on July 1, 2024, vacating the D.C. Circuit's judgment and remanding for further proceedings. Following Trump's reelection in November 2024, the Department of Justice moved under Federal Rule of Criminal Procedure 48(a) to dismiss the superseding indictment without prejudice pursuant to longstanding Office of Legal Counsel policy against prosecuting a sitting President, and Judge Chutkan granted the dismissal on November 25, 2024 (\href{https://clearinghouse.net/case/46788/}{Civil Rights Litigation Clearinghouse}).

\section{Judicial Votes}

The Court ruled \textbf{6--3} to vacate and remand, and \textbf{5--4} as to Part III-C of Chief Justice Roberts's opinion (barring evidentiary use of immune official acts) (\href{https://www.law.cornell.edu/supremecourt/text/23-939}{Cornell LII}):
\begin{itemize}
\item \textbf{Majority Opinion:} Chief Justice John G. Roberts, Jr., joined in full by Justices Clarence Thomas, Samuel A. Alito, Jr., Neil M. Gorsuch, and Brett M. Kavanaugh, and joined \emph{except as to Part III-C} by Justice Amy Coney Barrett.
\item \textbf{Concurring Opinion:} Justice Clarence Thomas (joining the Court's opinion in full and writing separately to question the validity of the Special Counsel's appointment under the Appointments Clause).
\item \textbf{Opinion Concurring in Part:} Justice Amy Coney Barrett (joining Parts I, II, III-A, III-B, IV, and V, while declining to join Part III-C's evidentiary bar and proposing a narrower two-step as-applied constitutional framework).
\item \textbf{Principal Dissent:} Justice Sonia Sotomayor, joined by Justices Elena Kagan and Ketanji Brown Jackson.
\item \textbf{Separate Dissent:} Justice Ketanji Brown Jackson (joining Justice Sotomayor's dissent in full and writing separately to critique the structural shift in presidential accountability).
\end{itemize}

\section{Holding}

\textbf{Judgment:} Vacated and remanded (91 F.4th 1173, vacated and remanded).

\textbf{Rule Announced:} Under the constitutional structure of separated powers, a former President's amenability to criminal prosecution turns on a tripartite framework (\href{https://www.law.cornell.edu/supremecourt/text/23-939}{603 U.S. 593}):
\begin{enumerate}
\item \textbf{Core Exclusive Powers (Absolute Immunity):} A former President is absolutely immune from criminal prosecution for conduct within his ``conclusive and preclusive'' constitutional authority under Article II, which Congress may not criminalize and courts may not examine.
\item \textbf{Remaining Official Acts (At Least Presumptive Immunity):} For all other official acts within the ``outer perimeter'' of presidential responsibility (conduct not ``manifestly or palpably beyond'' his authority), a former President is entitled to at least presumptive immunity, which the Government can rebut only by showing that applying a criminal prohibition to the act poses no ``dangers of intrusion on the authority and functions of the Executive Branch.''
\item \textbf{Unofficial Acts (No Immunity):} A former President enjoys no immunity from criminal prosecution for unofficial (private) acts.
\end{enumerate}

In distinguishing official from unofficial acts, courts may not inquire into the President's motives, nor may they deem an act unofficial merely because it allegedly violates a generally applicable criminal law. Further, a five-Justice majority (Part III-C) held that testimony or private records of the President or his advisers probing immune official acts may not be admitted as evidence at trial, even in a prosecution predicated solely on unofficial conduct. Finally, the Court rejected Trump's claim that the Impeachment Judgment Clause requires prior impeachment and Senate conviction before any criminal prosecution, and rejected the Government's claim that pretrial interlocutory review is unnecessary (\href{https://www.law.cornell.edu/supremecourt/text/23-939}{Cornell LII}).

\textbf{Holding vs. Dicta:}
\begin{itemize}
\item \textbf{Holdings:} (1) The tripartite immunity framework (absolute immunity for core exclusive powers, at least presumptive immunity for remaining official acts, no immunity for unofficial acts); (2) Trump is absolutely immune for his alleged discussions with Justice Department officials and threats to remove the Acting Attorney General, as investigative/prosecutorial direction and removal of principal executive officers fall within the President's conclusive and preclusive authority; (3) Trump's alleged discussions with the Vice President regarding the January 6 certification constitute official conduct entitled to at least presumptive immunity (remanded on whether the presumption is rebutted); (4) the Impeachment Judgment Clause does not make Senate conviction a prerequisite to criminal prosecution; (5) immunity rulings are immediately appealable before trial; and (6) prosecutors may not introduce immune official-act testimony or private records as evidence to prove unofficial-act charges (Part III-C).
\item \textbf{Dicta / Reserved Questions:} The Court expressly reserved whether immunity for non-core official acts is merely presumptive or actually absolute, noting that ``one case'' in more than ``two centuries does not afford enough experience'' to decide that question definitively. The Court also declined to decide whether Trump's interactions with state officials, private actors (the alternate elector scheme), or his January 6 Tweets and speech were official or unofficial, remanding those fact-bound classifications to the District Court. In Footnote 1, the Court noted that on remand the District Court should consider whether the \S 1512(c)(2) counts could proceed in light of \textit{Fischer v. United States}, 603 U.S. 480 (2024). Finally, Justice Thomas's concurrence questioning the Special Counsel's appointment was non-binding dicta, though it directly influenced Judge Aileen Cannon's subsequent July 15, 2024, dismissal of Trump's Florida classified-documents prosecution (\href{https://electionlawblog.org/2024/breaking-judge-cannon-says-appointment-of-special-counsel-was-unconstitutional-dismisses-trump-documents-case/}{Election Law Blog}).
\end{itemize}

\section{Analysis of Opinions}

\subsection{Majority Opinion (Chief Justice Roberts)}

Chief Justice Roberts anchored the Court's opinion in the structural imperatives of Article II and the separation of powers rather than an express constitutional immunity clause. Drawing on Justice Robert Jackson's concurring framework in \textit{Youngstown Sheet \& Tube Co. v. Sawyer}, 343 U.S. 579 (1952), Roberts reasoned that where the President exercises ``conclusive and preclusive'' constitutional authority---such as granting pardons (\textit{United States v. Klein}, 13 Wall. 128 (1872)), recognizing foreign governments (\textit{Zivotofsky v. Kerry}, 576 U.S. 1 (2015)), removing appointed executive officers (\textit{Myers v. United States}, 272 U.S. 52 (1926); \textit{Seila Law LLC v. CFPB}, 591 U.S. 197 (2020)), and directing Justice Department investigations and prosecutions---Congress is disabled from acting upon the subject and courts cannot adjudicate a criminal prosecution examining those acts (\href{https://www.law.cornell.edu/supremecourt/text/23-939}{Cornell LII}).

For official acts outside that exclusive core (where executive and legislative authority overlap), Roberts relied on Alexander Hamilton's call in \textit{The Federalist No. 70} for an ``energetic'' and ``vigorous'' Executive, as well as the Court's civil immunity precedent in \textit{Nixon v. Fitzgerald}, 457 U.S. 731 (1982), and executive-privilege decisions in \textit{United States v. Burr}, 25 F. Cas. 187 (C.C. Va. 1807), and \textit{United States v. Nixon}, 418 U.S. 683 (1974). Roberts reasoned that the threat of criminal trial and imprisonment is an even greater deterrent than civil damages and would chill ``bold and unhesitating action,'' risking ``an Executive Branch that cannibalizes itself, with each successive President free to prosecute his predecessors'' (\href{https://www.law.cornell.edu/supremecourt/text/23-939}{Cornell LII}). Balancing this against the compelling public interest in fair and effective law enforcement, the Court adopted at least presumptive immunity for all official acts within the outer perimeter, prohibiting judicial inquiry into presidential motives. Applying this framework, Roberts held Trump absolutely immune for his DOJ communications, presumptively immune for his interactions with Vice President Pence, and remanded the state-official, private-actor, and January 6 speech allegations. In Part III-C, Roberts (writing for five Justices) barred prosecutors from introducing testimony or private records probing immune official acts as evidence of unofficial crimes, reasoning that juries could otherwise indirectly judge immune official conduct, though Footnote 3 allowed prosecutors to cite the public record of an official act in a bribery case.

\subsection{Concurring Opinion (Justice Thomas)}

Justice Thomas joined the Court's opinion in full but wrote separately to question whether Special Counsel Jack Smith's office was ``established by Law'' as required by the Appointments Clause (U.S. Const. art. II, \S 2, cl. 2) (\href{https://www.law.cornell.edu/supremecourt/text/23-939}{Cornell LII}). Tracing the clause to the Founders' rejection of the British monarch's unilateral power to create offices, Thomas argued that none of the general statutory provisions cited by Attorney General Garland (28 U.S.C. \S\S 509, 510, 515, 533) clearly created the office of Special Counsel, unlike the lapsed Independent Counsel Act. Even if an office existed, Thomas urged lower courts to determine whether the Special Counsel was a principal officer requiring Senate confirmation or an inferior officer validly appointed by a Department Head before permitting the prosecution to proceed.

\subsection{Opinion Concurring in Part (Justice Barrett)}

Justice Barrett joined all of the majority opinion except Part III-C, offering a narrower doctrinal framing (\href{https://www.law.cornell.edu/supremecourt/text/23-939}{Cornell LII}; \href{https://reason.com/volokh/2024/07/05/what-exactly-did-justice-barrett-agree-with-chief-justice-roberts-on-about-presidential-immunity/}{Volokh Conspiracy}). She argued that ``immunity'' is simply shorthand for two propositions: (1) a President's right to bring an as-applied constitutional challenge to a criminal statute, and (2) the right to interlocutory appellate review before trial. Rather than leaving open whether non-core official acts carry absolute immunity, Barrett would have answered no and applied a two-step test: first, determine via ordinary statutory interpretation whether the criminal statute reaches the President's official conduct (without applying a clear-statement canon); second, if it does, ask whether prosecuting that conduct poses a danger of intrusion on Executive Branch authority. In Footnote 1, she stressed that not all exercises of the Take Care power are ``core,'' and in Footnote 2 she concluded that organizing fraudulent slates of electors is plainly private conduct because the President has zero constitutional role in state elector appointment. Crucially, Barrett rejected Part III-C's evidentiary bar, warning that blinding juries to an official act would hamstring prosecutions such as bribery (18 U.S.C. \S 201), where the jury must hear both the \textit{quid} and the \textit{quo}; ordinary evidentiary tools (Federal Rules of Evidence 403 and 105) suffice to prevent jury prejudice.

\subsection{Principal Dissent (Justice Sotomayor, joined by Justices Kagan and Jackson)}

Justice Sotomayor delivered a blistering dissent attacking the majority's official-acts immunity as ``atextual, ahistorical, and unjustifiable'' (\href{https://www.law.cornell.edu/supremecourt/text/23-939}{Cornell LII}). Textually, the Framers expressly granted legislative immunity in the Speech or Debate Clause (Art. I, \S 6, cl. 1) while omitting any presidential criminal immunity, and the Impeachment Judgment Clause explicitly contemplates subsequent criminal prosecution of former Presidents for official misconduct such as bribery. Historically, Alexander Hamilton (\textit{The Federalist No. 69}), Charles Pinckney, James Wilson, Justice Joseph Story, and modern executive practice---epitomized by President Ford's pardon of President Nixon and Nixon's acceptance---all confirmed that former Presidents are amenable to criminal law. On functionalist terms, Sotomayor argued that the majority inverted \textit{Fitzgerald}'s balancing test by ignoring the far weightier public interest in criminal law enforcement and the robust procedural safeguards (DOJ standards, grand juries, unanimous jury verdicts, beyond-a-reasonable-doubt proof) that prevent vexatious prosecutions. She further condemned the majority's expansion of ``core'' powers to sweep in the Take Care Clause and DOJ direction, its refusal to classify even concededly private acts in the indictment as unofficial, and Part III-C's ``nonsensical'' evidentiary bar. Famously warning that a President could order SEAL Team 6 to assassinate a political rival, organize a military coup, or take a bribe for a pardon and remain ``[i]mmune, immune, immune,'' Sotomayor concluded: ``With fear for our democracy, I dissent.''

\subsection{Separate Dissent (Justice Jackson)}

Justice Jackson joined Sotomayor's dissent in full and wrote separately to examine the theoretical shift from what she termed the ``individual accountability model''---under which the legislature defines crimes applicable to all citizens equally, subject to case-specific defenses---to a newly invented ``Presidential accountability model'' (\href{https://www.law.cornell.edu/supremecourt/text/23-939}{Cornell LII}; \href{https://www.lawfaremedia.org/article/the-supreme-court-s-presidential-immunity-decision}{Lawfare}). Because immunity means the law does not apply to the actor at all, Jackson argued that the majority revived the rejected monarchical maxim that ``the King can do no wrong,'' stripped Congress of its power to deter executive lawlessness, and installed the Supreme Court as the ultimate ad hoc gatekeeper deciding on a case-by-case basis which presidential acts are ``official'' and which prosecutions may proceed.

\section{Examples: Future Applications}

\begin{enumerate}
\item \textbf{Same-Side (Immune): Ordering the Attorney General to Drop a Criminal Investigation for Corrupt Political Reasons.} Suppose a sitting President directs the Attorney General to terminate a federal fraud investigation into a major campaign donor and threatens to fire the Attorney General if she refuses. After leaving office, the former President is indicted for obstruction of justice. Under \textit{Trump v. United States}, the former President is \textbf{absolutely immune}. Because directing DOJ investigative and prosecutorial decisions and removing principal executive officers fall within the President's ``conclusive and preclusive'' Article II authority, courts may neither criminalize the conduct nor inquire into the President's corrupt motives.
\item \textbf{Same-Side (Immune): Issuing a Pardon to a Co-Conspirator.} Suppose a President pardons a senior Cabinet official to prevent that official from testifying against the President in a congressional inquiry, and a subsequent administration charges the former President with witness tampering or obstruction. Under \textit{Trump v. United States}, the former President is \textbf{absolutely immune}. The pardon power (Art. II, \S 2, cl. 1) is a quintessential ``conclusive and preclusive'' power; Congress cannot regulate or criminalize its exercise, and courts cannot examine the President's motives for granting the pardon.
\item \textbf{Opposite-Side (Not Immune): Falsifying Personal Business Records Before or During Office.} Suppose a President, while in office, conspires with private accountants to submit fraudulent tax returns and falsified loan applications for a family-owned real estate business. After leaving office, the former President is indicted for bank fraud and tax evasion. Under \textit{Trump v. United States}, there is \textbf{no immunity}. Managing private family business finances is manifestly beyond any constitutional or statutory presidential responsibility and therefore constitutes purely unofficial conduct, just as in \textit{Clinton v. Jones}.
\item \textbf{Opposite-Side (Not Immune): Hiring a Private Party to Wiretap a Political Opponent's Campaign Headquarters.} Suppose a President running for reelection uses personal campaign funds to pay a private investigator to break into and wiretap the rival party's national headquarters, without involving any federal agency or invoking any national security authority. Under \textit{Trump v. United States}, there is \textbf{no immunity}. Conducting clandestine sabotage through private campaign operatives is an unofficial act undertaken as a candidate rather than an exercise of Article II authority.
\item \textbf{Fence-Sitter: Threatening to Withhold Congressionally Appropriated Foreign Aid to Pressure a Foreign Leader to Investigate a Domestic Rival.} Suppose a President pauses military aid appropriated by Congress for a foreign ally and tells the foreign leader in an official diplomatic phone call that release of the aid depends on announcing a corruption investigation into the President's domestic political rival. After leaving office, the former President is indicted under federal bribery or honest-services fraud statutes. On the one hand, conducting foreign diplomacy and communicating with foreign heads of state are squarely within the outer perimeter of official Article II responsibilities---and courts cannot examine the President's subjective motives to reclassify the call as unofficial---making the President at least \textbf{presumptively immune}. On the other hand, Congress holds concurrent authority over appropriations (Art. I, \S 9) and foreign commerce, and if the solicitation of a personal electoral benefit is viewed as a \textit{quid pro quo} outside exclusive executive power, a court might ask whether the Government can rebut the presumption (though Part III-C's evidentiary bar on probing the official diplomatic call would make proving the offense extraordinarily difficult).
\end{enumerate}

\section{Critique}

\textbf{Originalist and Textualist Critiques:} Scholars committed to original public meaning and textualism have criticized the majority's methodology from multiple angles. Conservative legal scholar Jack Goldsmith observed that the Court's sweeping characterization of exclusive presidential power over DOJ investigations and removals represents some of the most far-reaching pronouncements on executive power in Supreme Court history, resting heavily on functionalist fear of executive ``cannibalization'' rather than Founding-era text (\href{https://www.journals.uchicago.edu/doi/10.1086/735878}{Jack Goldsmith, \textit{The Supreme Court Review} (2025)}; \href{https://www.lawfaremedia.org/article/broad-reflections-on-trump-v.-united-states}{Lawfare}). Originalist scholar William Baude similarly noted that the Court's reasoning went well beyond any specific constitutional text or determinate tradition, relying instead on structural inferences and consequentialism. Professor Shalev Gad Roisman demonstrated in the \textit{University of Pennsylvania Law Review Online} that \textit{Trump v. United States} exposes the fatal limits of separation-of-powers formalism: criminal immunity disputes inherently arise in areas of overlapping congressional and presidential authority, yet the majority arbitrarily labeled DOJ direction and removal ``exclusive'' without originalist justification, and then abandoned categorical formalism altogether in favor of interest balancing for presumptive immunity (\href{https://pennlawreview.com/2025/02/03/trump-v-united-states-and-the-separation-of-powers/}{Penn Law Review Online}). Conversely, defenders such as Jonathan H. Adler argued that while immunity lacks an explicit textual clause, a narrower version---closer to Justice Barrett's concurrence---follows logically from the structural nature of Article II executive power.

\textbf{Progressive and Doctrinal Critiques:} Progressive and administrative-law scholars have condemned the decision as an unprecedented threat to democratic accountability and the rule of law. Writing in the \textit{Harvard Law Review Forum}, Professor Gillian E. Metzger (responding to Aziz Huq) argued that the majority engaged in ``faux balancing'' that distorted \textit{Nixon v. Fitzgerald} and \textit{Trump v. Vance}, sidelined Congress from defining presidential accountability, and undermined internal executive-branch legal guardrails by weakening the leverage of government lawyers who advise Presidents against unlawful acts (\href{https://harvardlawreview.org/forum/vol-138/disqualification-immunity-and-the-presidency/}{Harvard Law Review Forum}). Critics also highlight the stark methodological inconsistency between the Roberts Court's strict text-and-history originalism in individual-rights cases (\textit{Dobbs}, \textit{Bruen}) and its freewheeling, consequentialist structuralism in shielding executive power.

\textbf{Underlying Competing Values:} The decision reflects a profound collision between two foundational constitutional values: (1) \textit{executive energy, independence, and institutional stability}---ensuring that Presidents can make high-stakes decisions fearlessly without being paralyzed by the threat of partisan post-office prosecutions by political successors; and (2) \textit{the rule of law and democratic accountability}---ensuring that no official, least of all the chief magistrate entrusted to ``take Care that the Laws be faithfully executed,'' can abuse the awesome machinery of the state to subvert elections or commit crimes with impunity.

\section{Key Quotations}

\begin{quote}
``We conclude that under our constitutional structure of separated powers, the nature of Presidential power requires that a former President have some immunity from criminal prosecution for official acts during his tenure in office. At least with respect to the President's exercise of his core constitutional powers, this immunity must be absolute. As for his remaining official actions, he is also entitled to immunity.''
\end{quote}
\noindent --- Chief Justice John G. Roberts, Jr., opinion of the Court (\href{https://www.law.cornell.edu/supremecourt/text/23-939}{Cornell LII}; \href{https://www.courtlistener.com/opinion/10600038/trump-v-united-states/}{CourtListener}).

\begin{quote}
``The President enjoys no immunity for his unofficial acts, and not everything the President does is official. The President is not above the law. But Congress may not criminalize the President's conduct in carrying out the responsibilities of the Executive Branch under the Constitution.''
\end{quote}
\noindent --- Chief Justice John G. Roberts, Jr., opinion of the Court (\href{https://www.law.cornell.edu/supremecourt/text/23-939}{Cornell LII}).

\begin{quote}
``The Constitution does not require blinding juries to the circumstances surrounding conduct for which Presidents can be held liable.''
\end{quote}
\noindent --- Justice Amy Coney Barrett, concurring in part (\href{https://www.law.cornell.edu/supremecourt/text/23-939}{Cornell LII}).

\begin{quote}
``The relationship between the President and the people he serves has shifted irrevocably. In every use of official power, the President is now a king above the law.''
\end{quote}
\noindent --- Justice Sonia Sotomayor, dissenting (\href{https://www.law.cornell.edu/supremecourt/text/23-939}{Cornell LII}).

\begin{quote}
``Thus, being immune is not like having a defense under the law. Rather, it means that the law does not apply to the immunized person in the first place.''
\end{quote}
\noindent --- Justice Ketanji Brown Jackson, dissenting (\href{https://www.law.cornell.edu/supremecourt/text/23-939}{Cornell LII}).

\end{document}
